\begin{textsample}{POS Dim 2 – 2017\_c5 – Score -2.00 – 2017\_c5/enade\_2017\_mus\_ef85743bfe.txt}  \label{ex:f2_pos_394}
CAPÍTULO 5 DISTRIBUIÇÃO DOS CONCEITOS 5.1 PANORAMA NACIONAL DA DISTRIBUIÇÃO DOS CONCEITOS O Conceito Enade é um indicador de qualidade que avalia os cursos por intermédio dos desempenhos dos estudantes a partir dos resultados do Enade.
Os valores possíveis vão de 1 ( pior situação ) a 5 ( melhor situação ).
Os cursos com apenas um ou sem nenhum concluinte participante não obtêm o Conceito Enade, ficando Sem Conceito ( SC ).
Na Tabela 5.1, são apresentadas a quantidade e a distribuição de cursos de Música participantes do Enade/2017, por faixa de conceito e Grande Região.
Enfatiza -se, mais uma vez, que a diferença entre os cursos tabulados neste capítulo e no capítulo 2 corresponde aos cursos Sem Conceito ( SC ), ou seja àqueles, em princípio, sem alunos concluintes que houvessem participado da prova ou àqueles que tivessem tão somente um aluno concluinte18.
Observando -se os dados da Tabela 5.1, nota -se que, dos 107 cursos participantes, 39 ( 36,4 \% ) classificaram -se com conceito 2, o valor modal.
Esse foi também o conceito modal nas regiões Norte ( 100,0 \% ), Nordeste ( 55,0 \% ) e Sudeste ( 29,4 \% ).
Nas regiões Sul e Centro-Oeste, a moda foi o conceito 3, atribuído, respectivamente, a 31,8 \% e a 40,0 \% dos cursos da região.
O conceito 3 foi o segundo mais frequente em nível nacional ( 29,0 \% ), correspondendo a 31 cursos.
O conceito 4 foi o terceiro mais frequente ( 15,0 \%, correspondendo a 16 cursos ).
Houve, ainda, sete cursos ( 6,5 \% ) que receberam conceito 5, e oito cursos ( 7,5 \% ) que receberam conceito 1.
Dos 107 cursos de Música, seis ( 5,6 \% ) ficaram Sem Conceito ( SC ). 18 Esses cursos com somente um concluinte não foram considerados no capítulo 2.
A região Norte participou com quatro cursos ou 3,7 \% do total nacional.
Todos os quatro cursos receberam conceito 2, ou seja, 100,0 \% do total regional.
A região Nordeste participou com 20 cursos ou 18,7 \% do total nacional.
Nessa região, 11 cursos ( 55,0 \%, em termos regionais ) obtiveram conceito 3, o conceito modal para a região, como já comentado.
Os conceitos 1 e 4 foram igualmente atribuídos a um curso, cada um ( 5,0 \% para cada um ).
Ainda nessa região, sete cursos ( 35,0 \% ) foram avaliados com o conceito 4.
Nenhum curso ficou Sem Conceito ( SC ) ou com conceito 5.
Dos 51 cursos participantes da região Sudeste ( 47,7 \% do total nacional ), o conceito 2 foi o conceito modal, como já anteriormente assinalado, tendo sido obtido por 15 cursos ( 29,4 \% ).
O conceito 1 foi atribuído a seis cursos ( 11,8 \% ), o conceito 3, a 13 cursos ( 25,5 \% ).
Receberam os conceitos 4 e 5, respectivamente, oito cursos ( 15,7 \% ) e quatro cursos ( 7,8 \% ).
Nessa região, cinco cursos ( 9,8 \% ) ficaram Sem Conceito ( SC ).
Os 22 cursos da região Sul corresponderam a 20,6 \% do total nacional.
Como já indicado, houve predominância do conceito 3, atribuído a 31,8 \% dos cursos da região ( conceito modal ), ou seja, atribuído a sete dos 22 cursos participantes na região Sul.
O conceito 1 foi atribuído a um curso ( 4,5 \% ), o conceito 2, a seis cursos ( 27,3 \% ), o conceito 4, a cinco cursos ( 22,7 \% ), e o conceito 5, a três cursos ( 13,6 \% ).
Nenhum curso ficou Sem Conceito ( SC ) na região Sul.
Já dos dez cursos participantes na região Centro-Oeste ( 9,3 \% do total nacional ), quatro receberam conceito 3, o conceito modal ( 40,0 \% ), como já fora destacado.
A três cursos foi atribuído conceito 2 ( 30,0 \% ), e a dois cursos ( 20,0 \% ), conceito 4.
Os conceitos restantes, 1 e 5, não foram atribuídos a nenhum curso.
Nessa região, um curso ( 10,0 \% ) ficou Sem Conceito ( SC ).
No Gráfico 5.1, é apresentada a distribuição cumulativa do Conceito Enade, segundo a Grande Região.
Em um gráfico de distribuição cumulativa, em cada valor de abscissa ( eixo x ), o valor apresentado no eixo das ordenadas ( eixo y ) é o valor acumulado da variável.
Por exemplo, a região Nordeste ( linha verde ) apresenta 95 \% dos cursos com conceito 3 ou abaixo deste valor.
Por outro lado, a região Sul ( linha roxa ) apresenta um pouco mais de 63 \% dos cursos na mesma situação : com conceito 3 ou abaixo deste valor.
As regiões Nordeste ( linha verde ) e Centro-Oeste ( linha preta ) apresentam cursos com conceito 4 ou abaixo disso, já que a poligonal cumulativa alcança 100 \% para o conceito 4.
Nota -se que as regiões Norte ( linha azul ) e Centro-Oeste ( linha preta ) não apresentaram cursos com conceito 1, já que suas poligonais partem de 0 \% no conceito 1.
Vale ressaltar ainda que a poligonal cumulativa da região Norte alcançou a totalidade de cursos ( 100 \% ) no conceito 2, indicando que os cursos dessa região obtiveram apenas o conceito 2.
Em linhas gerais, num gráfico de distribuição cumulativa, poligonais mais à esquerda representam Grandes Regiões com pior distribuição, e poligonais mais à direita, Grandes Regiões com melhores conceitos.
A região Norte apresenta os cursos com conceitos mais baixos, seguido da região Nordeste, ao passo que a região Sul apresenta as distribuições com valores maiores.
Cursos em IES situadas nas regiões Sudeste ( linha vermelha ) e Centro-Oeste ( linha preta ) apresentam uma situação intermediária. 5.2 CONCEITOS POR CATEGORIA ADMINISTRATIVA, POR MODALIDADE DE ENSINO E GRANDE REGIÃO Na Tabela 5.2, é apresentada a distribuição dos cursos de Música participantes do Enade/2017, por Categoria Administrativa e Modalidade de ensino, de acordo com os conceitos por eles alcançados, segundo as Grandes Regiões brasileiras.
Dos 107 cursos participantes, 66 ( 61,7 \% ) eram ministrados em Instituições Públicas e 41 ( 38,3 \% ), em Instituições Privadas.
Quanto à Modalidade de ensino, 101 cursos eram presenciais ( 94,4 \% ), e seis ( 5,6 \% ), a distância.
De acordo com as informações da Tabela 5.2, em termos nacionais, entre Instituições Públicas e Privadas, observa -se que, dos sete cursos avaliados com conceito 5, seis eram oferecidos em IES Públicas e um, em IES Privada.
Dos 66 cursos participantes de IES Públicas, o conceito 2 foi o valor modal, atribuído a 23 cursos ( 34,8 \% ), como já comentado.
Entre os demais cursos de Instituições Públicas participantes, dois cursos ( 3,0 \% ) foram avaliados com conceito 1, outros 21 cursos ( 31,8 \% ) obtiveram conceito 3, 11 cursos ( 16,7 \% ) receberam conceito 4, e seis cursos ( 9,1 \% ) receberam conceito 5.
Nessa categoria, três cursos ( 4,5 \% ) ficaram Sem Conceito ( SC ).
Na rede privada, o conceito modal foi 2, atribuído a 16 cursos, o correspondente a 39,0 \% dos cursos da categoria.
Entre os demais cursos participantes, seis cursos ( 14,6 \% ) receberam conceito 1.
Entre os demais participantes, dez cursos ( 24,4 \% ) receberam conceito 3, cinco ( 12,2 \% ) receberam conceito 4, e apenas um ( 2,4 \% ), o conceito 5.
Nessa Categoria Administrativa, três cursos ( 7,3 \% ) ficaram Sem Conceito ( SC ).
Dos cursos participantes da modalidade presencial ( 94,4 \% do total de cursos ), o conceito 2 foi o valor modal, atribuído a 35 cursos ( 34,7 \% ).
Entre os demais cursos nessa modalidade de ensino, 14,9 \% foram avaliados com conceito 4, 29,7 \% obtiveram conceito 3, 6,9 \% receberam conceito 5, e outros 7,9 \% receberam conceito 1.
Nessa modalidade, seis cursos ( 5,9 \% ) ficaram Sem Conceito ( SC ).
Nos cursos a distância, o conceito modal foi também 2, atribuído a quatro cursos, correspondendo a 66,7 \% dos cursos da modalidade.
Dos demais cursos participantes, um deles 16,7 \% receberam conceito 3, e o outro ( 16,7 \% também ), conceito 4.
Nessa Modalidade, nenhum curso ficou Sem Conceito ( SC ) ou recebeu os conceitos 1 ou 5.
Na análise por região, observa -se que, na região Norte, que representou 3,7 \% do total de cursos participantes em nível nacional, não houve participação de cursos de Instituições Privadas de ensino.
As Instituições Públicas participaram com quatro cursos, todos presenciais, que obtiveram o conceito 2.
A região Nordeste, representando 18,7 \% do total de cursos, contou com dois dos 20 cursos participantes em Instituições Privadas.
Ambos obtiveram conceito 2.
As Instituições Públicas participaram com 18 cursos.
Desses, nove ( 50,0 \% ) obtiveram conceito 2, o valor modal.
Na sequência, sete cursos obtiveram conceito 3 ( 38,9 \% ).
Os conceitos 1 e 4 foram atribuídos a um curso, cada um ( 5,6 \% ).
Nenhum curso ficou Sem Conceito ( SC ) ou recebeu conceito 5.
Todos os 20 cursos na região foram presenciais.
A região Sudeste, que representou quase a metade do total de cursos participantes em nível nacional ( 47,7 \% ), apresentou mais da metade deles vinculados à rede privada, 54,9 \%, correspondendo a 28 dos 51 cursos participantes.
Nessa categoria e região, o conceito modal foi 2, atribuído a 12 cursos ( 42,9 \% ).
Dois cursos ( 7,1 \% ) ficaram Sem Conceito ( SC ).
Os demais foram avaliados com conceito 1 ( cinco cursos, 17,9 \% ), conceito 3 ( seis cursos, 21,4 \% ), conceito 4 ( dois cursos, 7,1 \% ) e conceito 5 ( um curso, 3,6 \% ).
Entre os 23 cursos oferecidos por Instituições Públicas na região Sudeste ( 45,1 \% ), o conceito modal foi 3, atribuído a sete cursos ( 30,4 \% ).
Os demais cursos receberam o conceito 1 ( um curso, 4,3 \% ), 2 ( três cursos, 13,0 \% ), 4 ( seis cursos, 26,1 \% ) e 5 ( três cursos, 13,0 \% ).
Na região Sudeste, nessa categoria, três cursos ( 13,0 \% ) ficaram Sem Conceito ( SC ).
Também na região Sudeste, a maioria dos cursos foram presenciais : 90,2 \%.
Para esses, a moda foi atribuída aos conceitos 2 e 3, cada um com 12 cursos, correspondendo a 26,1 \% para cada um.
Os demais cursos presenciais da região foram alocados aos conceitos 1 ( 13,0 \% ), 4 ( 15,2 \% ) e 5 ( 8,7 \% ).
Dos cursos presenciais, cinco ( 10,9 \% ) ficaram Sem Conceito ( SC ).
Para os cinco cursos a distância ( 9,8 \% ), a moda também foi o conceito 2, alcançado por três cursos, 60,0 \%.
Os demais cursos a distância da região foram alocados aos conceitos 3 e 4, cada qual com um curso ( 20,0 \% para cada um ).
Nenhum curso a distância da região recebeu os conceitos 1 ou 5, e nenhum ficou Sem Conceito ( SC ).
As Instituições Privadas concentraram dez dos 22 cursos participantes da região Sul, 45,5 \% do total regional.
Desses, quatro ( 40,0 \% ) obtiveram conceito 3, o conceito modal.
Nessa combinação de Categoria Administrativa e Grande Região, os conceitos 2 e 4 foram atribuídos a dois ( 20,0 \% ) e a três ( 30,0 \% ) cursos, respectivamente.
Um curso recebeu conceito 1 ( 10,0 \% ).
Nenhum recebeu conceito 5, assim como nenhum ficou Sem Conceito ( SC ).
As Instituições Públicas na região Sul participaram com 12 cursos, 54,5 \% dos cursos da região.
O conceito modal foi 2, atribuído a quatro cursos ( 33,3 \% ).
Aos demais cursos foram atribuídos os conceitos 3 ( três cursos, 25,0 \% ) 4 ( dois cursos, 16,7 \% ) e 5 ( três cursos, 25,0 \% ).
Nenhum curso da região Sul oferecido por IES Pública ficou Sem Conceito ( SC ) ou recebeu conceito 1.
Todos os 22 cursos da região Sul foram presenciais.
Na região Centro-Oeste, apenas um dos dez cursos participantes era de Instituições Privadas ( 10,0 \% em termos regionais ).
Esse único curso ficou Sem Conceito ( SC ).
Quanto aos nove cursos oferecidos por Instituições Públicas na região Centro-Oeste ( 44,4 \% do total regional ), quatro cursos apresentaram conceito 3, conceito modal.
Três cursos ( 33,3 \% ) obtiveram conceito 2, e dois cursos ( 22,2 \% ) obtiveram conceito 4.
Nenhum curso ficou Sem Conceito ( SC ) ou com conceitos 1 ou 5.
Nessa região, 90,0 \% dos cursos foram presenciais.
Para os cursos presenciais, a moda foi o conceito 3 com 44,4 \% dos cursos, correspondendo a 4 cursos.
Os outros cursos foram alocados igualmente aos conceitos 3 e 4, dois cursos em cada conceito, correspondendo a 22,2 \% cada.
Um curso ( 11,1 \% ) ficou Sem Conceito ( SC ).
O único curso na Modalidade a Distância foi alocado ao conceito 2.
No Gráfico 5.2, é apresentada a distribuição cumulativa do Conceito Enade segundo a Categoria Administrativa da IES.
Os cursos em IES Públicas ( linha azul ) apresentam uma poligonal mais à direita do que os em IES Privadas ( linha verde ), em que 57,9 \% dos cursos desta IES obtiveram conceito 2 ou abaixo deste valor e os daquela IES, 39,7 \%, indicando uma distribuição de Conceitos Enade com valores maiores. 5.3 CONCEITOS POR ORGANIZAÇÃO ACADÊMICA E GRANDE REGIÃO Na Tabela 5.3, encontra -se a distribuição dos conceitos atribuídos aos cursos participantes do Enade/2017 na Área de Música, por Organização Acadêmica, segundo as Grandes Regiões brasileiras.
Dos 107 cursos de Música participantes, 78 eram oferecidos em Universidades ; 12, em Centros Universitários ; 14, em Faculdades ; e três, em CEFET/IFET.
Essa distribuição corresponde a, respectivamente, 72,9 \%, 11,2 \% e 13,1 \% e 2,8 \% dos cursos.
De acordo com os dados apresentados, de todos os sete cursos avaliados com conceito 5, seis eram vinculados a Universidades.
Os cursos vinculados a esse tipo de Organização Acadêmica tiveram o conceito 2 como conceito modal, atribuído a 26 cursos ( 33,3 \% ).
Os demais cursos vinculados a Universidades receberam os conceitos 1 ( cinco cursos, 6,4 \% ), 3 ( 25 cursos, 32,1 \% ), 4 ( 13 cursos, 16,7 \% ) e conceito 5 ( seis cursos, 7,7 \%, como já mencionado ).
Ficaram Sem Conceito ( SC ) três cursos ( 3,8 \% ).
Dentre os cursos em Centros Universitários, o conceito modal foi também 2, atribuído a quatro cursos ( 33,3 \% ).
Um curso vinculado a esse tipo de Organização Acadêmica ( 8,3 \% ) ficou Sem Conceito ( SC ).
Os outros cursos receberam os conceitos 1 ( dois cursos, 16,7 \% ), 3 ( três cursos, 25,0 \% ), e 4 ( dois cursos,16,7 \% ).
Nenhum curso recebeu conceito 5.
Dos 14 cursos mantidos por Faculdades, sete ( 50,0 \% ) receberam conceito 2, o conceito modal.
O segundo conceito mais frequentemente atribuído aos cursos oferecidos por essa categoria de IES foi o conceito 3, recebido por dois cursos ( 14,3 \% ).
Na sequência, um curso ( 7,1 \% ) obteve conceito 4, e outro curso ( 7,1 \% ), conceito 5.
Dois cursos ( 14,3 \% ) ficaram Sem Conceito ( SC ) e um curso ( 7,1 \% ) recebeu conceito 1.
Dos três cursos mantidos por CEFET/IFET, dois ( 66,7 \% ) receberam conceito modal 2, e um curso ( 33,3 \% ) recebeu conceito 3.
Considerando -se separadamente as regiões brasileiras, verifica -se que, na região Norte, todos os quatro cursos participantes eram de Universidades e obtiveram o conceito 2.
Na região Nordeste, as Universidades participaram com 18 dos 20 cursos da Área de Música oferecidos na região ( 90,0 \% ), dos quais nenhum ficou Sem Conceito ( SC ) ou recebeu o conceito 5.
Dez cursos obtiveram o conceito modal 2 ( 55,6 \% ).
Um curso ( 5,6 \% ) obteve conceito 1, seis cursos ( 33,3 \% ) receberam conceito 3, e um curso ( 5,6 \% ) obteve conceito 4.
Os Centros Universitários e Faculdades não contaram com cursos participantes na região Nordeste.
Já os CEFET/IFET participaram com dois cursos, dos quais um recebeu o conceito 2 e o outro o conceito 3.
Na região Sudeste, as Universidades concentraram 31 dos 51 cursos de Música da região ( 60,8 \% ).
Dentre os cursos em Universidades, nessa região, o conceito modal foi 3, atribuído a dez cursos ( 32,3 \% ).
Os demais cursos receberam os conceitos 1 ( três cursos, 9,7 \% ), 2 ( cinco cursos, 16,1 \% ), 4 ( sete cursos, 22,6 \% ) e 5 ( três cursos, 9,7 \% ).
Três cursos ( 9,7 \% ) ficaram Sem Conceito ( SC ).
Os Centros Universitários participaram com dez cursos na região Sudeste ( 19,6 \% ), dos quais quatro ( 40,0 \% ) obtiveram o conceito modal 2.
O conceito 1 foi atribuído a dois cursos ( 20,0 \% ), o conceito 2, a dois cursos também ( 20,0 \% ), e o conceito 4, a um curso ( 10,0 \% ).
Nenhum curso obteve conceito 5, e um curso ( 10,0 \% ) ficou Sem Conceito ( SC ) nessa combinação de organização acadêmica e região.
As Faculdades foram representadas também por dez cursos na região Sudeste ( 19,6 \% ), dos quais um ( 10,0 \% ) ficou Sem Conceito ( SC ).
Seis cursos ( 60,0 \% ) obtiveram conceito 2.
Os conceitos 1, 3 e 5 foram atribuídos a um curso ( 10,0 \% ), cada um.
Nenhum curso recebeu conceito 4.
Os CEFET/IFET não participaram com nenhum curso na região.
Dos 22 cursos da região Sul, 17 ( 77,3 \% ) eram de Universidades, para os quais os conceitos 2 e 3 foram igualmente atribuídos a cinco cursos ( 29,4 \% ), cada um, conceitos modais.
Os conceitos 4 e 5 foram também igualmente distribuídos por três cursos ( 17,6 \% ), cada um, e um curso ( 5,9 \% ) recebeu conceito 1.
Nesse tipo de organização, nenhum curso ficou Sem Conceito ( SC ).
Dos dois cursos participantes de Centros Universitários da região Sul ( 9,1 \% ), um recebeu conceito 3, e outro, conceito 4.
Quanto aos três cursos vinculados a Faculdades na região Sul ( 13,6 \% ), os conceitos 2, 3 e 4 foram atribuídos a um curso, cada um.
Nessa região, não foram avaliados cursos em CEFET/IFET.
Na região Centro-Oeste, oito dos dez cursos eram de Universidades ( 80,0 \% ).
Nesse tipo de organização, nenhum curso ficou Sem Conceito ( SC ), assim como nenhum recebeu conceitos 1 ou 5.
Quatro cursos receberam o conceito modal 3 ( 50,0 \% ).
Os conceitos 2 e 4 foram atribuídos a dois cursos, cada um ( 25,0 \% ).
Não houve participação de cursos em Centros Universitários na região Centro-Oeste.
O único curso oferecido em Faculdades ficou Sem Conceito ( SC ).
Um único curso foi avaliado em CEFET/IFET, tendo recebido conceito 2.
No Gráfico 5.3, é apresentada a distribuição cumulativa do Conceito Enade segundo a Organização Acadêmica da IES.
Os cursos em Universidades ( linha azul ) apresentam uma poligonal mais à direita do que as demais ( 41,3 \% dos cursos avaliados com conceito 2 ou abaixo ), denotando melhor distribuição de Conceitos Enade.
No outro extremo, encontram -se os cursos oferecidos em CEFET/IFET ( linha roxa ), com a poligonal mais à esquerda e, consequentemente, uma distribuição de Conceitos Enade com menores valores.
Vale ressaltar que a poligonal dos cursos ofertados em CEFET/IFET ( linha roxa ) parte de 0 \% no conceito 1 e alcança 100 \% no conceito 3, indicando que os cursos dessa Organização Acadêmica não obtiveram conceitos 1, 4 e 5.
Já os cursos em Centros Universitários ( linha verde ) estão todos com conceito 4 ou abaixo, pois a poligonal que os representa alcança 100 \% no conceito 4.
Os IES de Faculdades ( linha vermelha ) apresentam uma situação intermediária.

% matched lemmas: 
\end{textsample}
